Controlling video and audio generation requires diverse modalities, from depth and pose to camera trajectories and audio transformations, yet existing approaches either train a single monolithic model for a fixed set of controls or introduce costly architectural changes for each new modality.
We introduce AVControl, a lightweight, extendable framework built on LTX-2~\cite{hacohen2026ltx2}, a joint audio-visual foundation model, where each control modality is trained as a separate LoRA on a \emph{parallel canvas} that provides the reference signal as additional tokens in the attention layers, requiring no architectural changes beyond the LoRA adapters themselves.
We show that simply extending image-based in-context methods to video fails for structural control, and that our parallel canvas approach resolves this.
On the VACE Benchmark~\cite{jiang2025vace}, we outperform all evaluated baselines on depth- and pose-guided generation, inpainting, and outpainting, and show competitive results on camera control and audio-visual benchmarks.
Our framework supports a diverse set of independently trained modalities: spatially-aligned controls such as depth, pose, and edges, camera trajectory with intrinsics, sparse motion control, video editing, and, to our knowledge, the first modular audio-visual controls for a joint generation model.
Our method is both compute- and data-efficient: each modality requires only a small dataset and converges within a few hundred to a few thousand training steps, a fraction of the budget of monolithic alternatives.
We publicly release our code and trained LoRA checkpoints.
